% Respuesta a la pregunta 5
\item
Las bases de datos NoSQL (\textit{Not Only SQL}) son sistemas de almacenamiento no relacionales diseñados para escalabilidad horizontal, flexibilidad de esquema y alto rendimiento en casos de uso específicos.

\textbf{3 Ventajas frente a bases relacionales:}
\begin{enumerate}
	\item \textbf{Escalabilidad horizontal nativa:} Sharding y replicación automáticos (MongoDB, Cassandra) sin rediseño de esquema.
	\item \textbf{Flexibilidad de esquema (schema-less):} Documentos JSON variables, ideal para datos semi-estructurados y desarrollo ágil.
	\item \textbf{Rendimiento en escrituras masivas:} Modelos como LSM-trees (Cassandra, RocksDB) optimizados para alta ingestión de datos.
\end{enumerate}

\textbf{3 Desventajas frente a bases relacionales:}
\begin{enumerate}
	\item \textbf{Consistencia eventual:} Modelo BASE vs ACID; transacciones distribuidas complejas o inexistentes.
	\item \textbf{Falta de estándar de consulta:} Cada NoSQL tiene su API/QL (MQL, CQL, Cypher, Redis commands), curva de aprendizaje por tecnología.
	\item \textbf{Integridad referencial limitada:} Sin foreign keys, cascadas, ni constraints declarativas; la lógica recae en la aplicación.
\end{enumerate}
